\documentclass[10pt,a4paper]{amsart}
\usepackage{amsmath,amssymb,amsthm,graphicx}
% Retained minimal preamble for the Stacks Project chapter more-algebra.tex (More on Algebra).
% The chapter file pinned-original-source.tex carries only \input{preamble} on its first line, so no
% preamble text lives inside the chapter itself.  The Stacks Project uses one single shared preamble for
% every chapter; the snapshot stages it identically as sources/stacks-chow/preamble.tex and
% sources/stacks-etale-cohomology/preamble.tex.  The block below is copied mechanically, line for line,
% from sources/stacks-chow/preamble.tex lines 14-18 (the Xy-pic setup) and 154-217 (the native theorem
% declarations together with the mathematical macros).  Nothing was added, removed or redefined; the
% wrapper class and the amsmath/amssymb/amsthm/graphicx base packages are supplied by the assembler.
% Chapter: more-algebra.tex -- More on Algebra.
% For commutative diagrams we use Xy-pic
\usepackage[all]{xy}

% We use 2cell for 2-commutative diagrams.
\xyoption{2cell}

% Theorem environments.
%
\theoremstyle{plain}
\newtheorem{theorem}[subsection]{Theorem}
\newtheorem{proposition}[subsection]{Proposition}
\newtheorem{lemma}[subsection]{Lemma}

\theoremstyle{definition}
\newtheorem{definition}[subsection]{Definition}
\newtheorem{example}[subsection]{Example}
\newtheorem{exercise}[subsection]{Exercise}
\newtheorem{situation}[subsection]{Situation}

\theoremstyle{remark}
\newtheorem{remark}[subsection]{Remark}
\newtheorem{remarks}[subsection]{Remarks}

\numberwithin{equation}{subsection}

% Macros
%
\def\lim{\mathop{\mathrm{lim}}\nolimits}
\def\colim{\mathop{\mathrm{colim}}\nolimits}
\def\Spec{\mathop{\mathrm{Spec}}}
\def\Hom{\mathop{\mathrm{Hom}}\nolimits}
\def\Ext{\mathop{\mathrm{Ext}}\nolimits}
\def\SheafHom{\mathop{\mathcal{H}\!\mathit{om}}\nolimits}
\def\SheafExt{\mathop{\mathcal{E}\!\mathit{xt}}\nolimits}
\def\Sch{\mathit{Sch}}
\def\Mor{\mathop{\mathrm{Mor}}\nolimits}
\def\Ob{\mathop{\mathrm{Ob}}\nolimits}
\def\Sh{\mathop{\mathit{Sh}}\nolimits}
\def\NL{\mathop{N\!L}\nolimits}
\def\CH{\mathop{\mathrm{CH}}\nolimits}
\def\proetale{{pro\text{-}\acute{e}tale}}
\def\etale{{\acute{e}tale}}
\def\QCoh{\mathit{QCoh}}
\def\Ker{\mathop{\mathrm{Ker}}}
\def\Im{\mathop{\mathrm{Im}}}
\def\Coker{\mathop{\mathrm{Coker}}}
\def\Coim{\mathop{\mathrm{Coim}}}

% Boxtimes
%
\DeclareMathSymbol{\boxtimes}{\mathbin}{AMSa}{"02}

%
% Macros for moduli stacks/spaces
%
\def\QCohstack{\mathcal{QC}\!\mathit{oh}}
\def\Cohstack{\mathcal{C}\!\mathit{oh}}
\def\Spacesstack{\mathcal{S}\!\mathit{paces}}
\def\Quotfunctor{\mathrm{Quot}}
\def\Hilbfunctor{\mathrm{Hilb}}
\def\Curvesstack{\mathcal{C}\!\mathit{urves}}
\def\Polarizedstack{\mathcal{P}\!\mathit{olarized}}
\def\Complexesstack{\mathcal{C}\!\mathit{omplexes}}
% \Pic is the operator that assigns to X its picard group, usage \Pic(X)
% \Picardstack_{X/B} denotes the Picard stack of X over B
% \Picardfunctor_{X/B} denotes the Picard functor of X over B
\def\Pic{\mathop{\mathrm{Pic}}\nolimits}
\def\Picardstack{\mathcal{P}\!\mathit{ic}}
\def\Picardfunctor{\mathrm{Pic}}
\def\Deformationcategory{\mathcal{D}\!\mathit{ef}}
\usepackage{cleveref}
\crefname{originallemma}{original lemma}{original lemmas}
\Crefname{originallemma}{Original lemma}{Original lemmas}
\makeatletter
\newlabel{algebra-lemma-hom-from-finitely-presented}{{0583}{source}{}{sourcepointer-1}{}}
\newlabel{algebra-lemma-hom-from-finitely-presented@cref}{{[originallemma][0][]0583}{[1][0][]source}{}{}{}}
\makeatother
\title{A tuple becomes a direct summand inclusion after inverting some element of the complement of a point exactly when its image generates an independent family of the fibre}
\author{The Stacks Project Authors}
\date{stacks-project:more-algebra.tex:local-snapshot-20261012; GNU FDL 1.2 or later}
\begin{document}
\maketitle
\setcounter{subsection}{1}
\noindent\textit{Editorial note.} All delivered mathematical text is the author's own, copied from the pinned Stacks Project chapter file more-algebra.tex; no statement or proof step was rewritten, reordered or removed. The counted statement is the labeled lemma environment at source lines 38730-38746 (canonical tag 0GV9) and the counted proof is the proof environment at source lines 38748-38786, which immediately follows it and is its own proof.\par\medskip
\noindent\textit{Editorial note.} The statement is formulated inside the situation of the section, which is why it begins by referring to Situation situation-splitting; that label is carried as a navigation stub because the preamble of this card is the shared chapter preamble and the situation environment itself is not counted.\par\medskip
\noindent\textit{Editorial note.} The proof contains one footnote, the Kronecker delta explaining the notation delta\_ij, at source line 38761. It is part of the author's own text inside the counted span and is retained.\par\medskip
\noindent\textit{Editorial note.} The chapter file pinned-original-source.tex carries only an input of the shared preamble on its first line, so no preamble text lives inside the chapter. The retained file sources/stacks-more-algebra/retained-preamble.tex is a mechanical, line-for-line copy of the project's single shared preamble as staged in the snapshot (Xy-pic setup plus the native theorem declarations and mathematical macros). Nothing was added, removed or redefined, and the excerpt serves only as preamble, never as a span coordinate.\par\medskip
\noindent\textit{Editorial note.} The source URL is the source-level Stacks Project URL, as required when more than one card is produced from one source; the per-claim tag is carried by claim\_key and by this note. The counted statement and proof cite no bibliography entry, so this card needs no bibliography.\par\medskip
\noindent\textit{Allowed standard prerequisites (editorial).}\begin{itemize}
\item A family of elements of a module is the inclusion of a direct summand after localizing at f exactly when there are dual maps whose composite matrix has determinant mapping to a unit of the localization. used at proof lines 38762-38772, where the composite of the tuple with the dual family has determinant f mapping to 1 in kappa(x), which gives the splitting after inverting f. Produces the direct summand condition from the independence of the images in V(x).
\item Localization commutes with Hom from a finitely presented module, so maps defined after inverting f can be cleared by multiplying with a power of f. used at proof lines 38778-38783, where the maps phi\_i over R\_f are replaced by maps phi'\_i over R with phi'\_i(s\_j) = f\textasciicircum{}n delta\_ij. Provides the converse direction: the dual family over the point comes from an actual family of maps over R after multiplying by a power of f.
\item The perpendicular B(x) of the image of the evaluation map consists of the linear forms vanishing on the images of all R-linear maps, and the quotient V(x) = M(x)/B(x) is the fibre that detects independence. used at proof lines 38749-38756. Fixes the objects in terms of which both directions of the equivalence are formulated.
\end{itemize}\medskip
\ifdefined\hypertarget\hypertarget{sourcepointer-1}{}\fi
\noindent\textit{Source reference.} Original lemma 0583: sources/stacks-chow/tags.txt, source lines 6766--6766. The reference supplies the standard compatibility of Hom from a finitely presented module with localization, used at proof lines 38780-38782 to clear the denominators f in the dual maps; the argument that consumes it, namely passing from the maps over R\_f to the maps phi'\_i over R and reading off independence, is inside the counted span.\par\medskip
\begin{situation}
\label{situation-splitting}
Here $R$ is a ring and $M$ is a finitely presented $R$-module.
Denote $\Omega \subset \Spec(R)$ the set of closed points
with the induced topology. For $x \in \Omega$ denote
$M(x) = M/xM$ the fibre of $M$ at $x$. This is a finite
dimensional vector space over the residue field $\kappa(x)$ at $x$.
Given $s \in M$ we denote $s(x)$ the image of $s$ in $M(x)$.
\end{situation}

\begin{lemma}
\label{lemma-which-elements-split}
In Situation \ref{situation-splitting} let $x \in \Omega$. There exists a
canonical short exact sequence
$$
0 \to B(x) \to M(x) \to V(x) \to 0
$$
of $\kappa(x)$-vector spaces which the following
property: for $s_1, \ldots, s_r \in M$ the following are equivalent
\begin{enumerate}
\item there exists an $f \in R$, $f \not \in x$ such that
the map $s_1, \ldots, s_r : R^{\oplus r} \to M$ becomes the
inclusion of a direct summand after inverting $f$, and
\item $s_1(x), \ldots, s_r(x)$ map to linearly independent
elements of $V(x)$.
\end{enumerate}
\end{lemma}

\begin{proof}
Define $B(x) \subset M(x)$ as the perpendicular of the image of the map
$$
\Hom_R(M, R) \to \Hom_{\kappa(x)}(M(x), \kappa(x))
$$
and set $V(x) = M(x)/B(x)$. Then any $R$-linear map $\varphi : M \to R$
induces a map $\overline{\varphi} : V(x) \to \kappa(x)$ and conversely
any $\kappa(x)$-linear map $\lambda : V(x) \to \kappa(x)$ is equal
to $\overline{\varphi}$ for some $\varphi$. Let $s_1, \ldots, s_r \in M$.

\medskip\noindent
Suppose $s_1, \ldots, s_r$ map to linearly independent elements of
$V(x)$. Then we can find $\varphi_1, \ldots, \varphi_r \in \Hom_R(M, R)$
such that $\varphi_i(s_j)$ maps to $\delta_{ij}$\footnote{Kronecker delta.}
in $\kappa(x)$. Hence the matrix of the composition
$$
R^{\oplus r}
\xrightarrow{s_1, \ldots, s_r}
M
\xrightarrow{\varphi_1, \ldots, \varphi_r}
R^{\oplus r}
$$
has a determinant $f \in R$ which maps to $1$ in $\kappa(x)$
Clearly, this implies that $s_1, \ldots, s_r : R^{\oplus r} \to M$
is the inclusion of a direct summand after inverting $f$.

\medskip\noindent
Conversely, suppose that we have an $f \in R$, $f \not \in x$
such that $s_1, \ldots, s_r : R^{\oplus r} \to M$
is the inclusion of a direct summand after inverting $f$.
Hence we can find $R_f$-linear maps $\varphi_i : M_f \to R_f$
such that $\varphi_i(s_j) = \delta_{ij} \in R_f$.
Since $\Hom_R(M, R)_f = \Hom_{R_f}(M_f, R_f)$ by
Algebra, Lemma \ref{algebra-lemma-hom-from-finitely-presented}
we conclude that we can find $n \geq 0$ and $\varphi'_i \in \Hom_R(M, R)$
such that $\varphi'_i(s_j) = f^n\delta_{ij} \in R$.
It follows that $s_1, \ldots, s_r$ map to linearly independent
elements of $V(x)$ as $\overline{\varphi}'_i(s_j) = f^n\delta_{ij}$.
\end{proof}

\end{document}
